% Generated draft: requires mathematical and source review.
\documentclass[11pt]{article}
\usepackage[margin=25mm]{geometry}
\usepackage{fontspec}
\setmainfont{Latin Modern Roman}
\usepackage{amsmath,amssymb,mathtools,unicode-math}
\setmathfont{Latin Modern Math}
\usepackage{enumitem,needspace,xurl,hyperref}
\setlength{\parindent}{0pt}
\setlength{\parskip}{5pt}
\setlength{\emergencystretch}{3em}
\newcommand{\N}{\mathbb{N}}
\newcommand{\Z}{\mathbb{Z}}
\newcommand{\Q}{\mathbb{Q}}
\newcommand{\R}{\mathbb{R}}
\newcommand{\C}{\mathbb{C}}
\newcommand{\F}{\mathbb{F}}
\newcommand{\A}{\mathbb{A}}
\newcommand{\PP}{\mathbb{P}}
\newcommand{\E}{\mathbb{E}}
\newcommand{\eps}{\varepsilon}
\newcommand{\dd}{\,\mathrm{d}}
\begin{document}
\section*{UnsolvedMath ID: 9400100}
\textbf{Source problem code:} catalogue entry
\subsection{Definitions and mathematical statement}
For a real number $x$, fix the irrationality-exponent convention
\[
\mu(x)=\sup\left\{\nu>0:\ 0<\left|x-\frac ab\right|<b^{-\nu}
\text{ for infinitely many coprime }a\in\Z,\ b\in\N_{\geq1}\right\}.
\]
The exclusion of zero error prevents repeated exact representations. This convention extends to rational numbers with value one; for irrationals the exponent is at least two. Some sources reserve the name irrationality measure for irrational inputs, so the convention matters for $\gamma$.

Let $\pi$ be the usual circle constant and let
\[
\gamma=\lim_{N\to\infty}\left(\sum_{j=1}^N\frac1j-\log N\right)
\]
be the Euler--Mascheroni constant. The concrete request is to determine the exact values of $\mu(\pi)$ and $\mu(\gamma)$, under a specified convention, and more broadly to understand rational approximation of nonquadratic irrationals. The summary reports irrationality of $\gamma$ itself unresolved, so it must not be assumed in defining the second target. Proving a value appropriate only to irrational inputs would require addressing that preliminary issue.

The nonquadratic class is not synonymous with transcendental numbers: it includes algebraic irrationals of degree at least three. The supplied summary invokes Roth's theorem for their exponent two and distinguishes that known class from specific transcendental or suspected-transcendental constants. No universal common exponent for every nonquadratic irrational is proposed. The exact-value question is also stronger than obtaining any finite upper bound for $\mu(\pi)$.
\subsection{Short English statement}
Determine exact rational-approximation exponents for $\pi$ and the Euler--Mascheroni constant. The latter may not yet be assumed irrational, and higher-degree algebraic irrationals are a separate known-exponent class.
\subsection{Sources}
\begin{itemize}
\item Local UnsolvedMath registry, permanent ID 9400100, exact source code catalogue entry; no proposer or year supplied. Registry status open, classification OPEN-TRIAGE.
\item Wikipedia contributors, \emph{List of unsolved problems in mathematics}, Diophantine approximation and transcendental number theory, bullet 100, wikitext line 126, pinned revision 1366636767, \url{https://en.wikipedia.org/w/index.php?title=List_of_unsolved_problems_in_mathematics&oldid=1366636767#Number_theory}. Supplied MediaWiki extraction accessed July 29, 2026; staged pending semantic deduplication and primary-source review.
\item The summary names Lambert for irrationality of $\pi$ and Roth for algebraic-irrational exponent two. It reports unknown exact measures and unresolved irrationality of $\gamma$, and gives only an approximate upper-bound description $\mu(\pi)\leq7.103\ldots$ without authors, date or precise theorem. That decimal is retained as supplied historical/progress information, not asserted to be the current sharpest bound.
\item No primary exact-measure citation is supplied. Provenance flags current primary-literature status and redistribution terms NEEDS\_REVIEW; the summary refers generally to Wikipedia's ``Irrationality measure''.
\end{itemize}
\end{document}
